% =========================================================
% Continuity --- Notes
% Chapter: continuity
% Section: continuity
% Volume III - Analysis
% =========================================================

\subsection*{Continuity}
\label{sec:continuity}

% ---------------------------------------------------------
% Toolkit
% ---------------------------------------------------------
% ---------------------------------------------------------
% Definition --- Relative Neighborhood
% ---------------------------------------------------------

\begin{definitionbox}{Definition (Relative Neighborhood)}
\begin{definition}[Relative Neighborhood]
\label{def:relative-neighborhood}
Let $E \subseteq \mathbb{R}$ and $a \in E$. A set $U \subseteq \mathbb{R}$ is a neighborhood of $a$ in $E$ if and only if there exists $\delta > 0$ such that $U = E \cap (a-\delta, a+\delta)$. Equivalently, every neighborhood of $a$ in $E$ is a set of the form $U_E(a) := E \cap (a-\delta, a+\delta)$ for some $\delta > 0$.
\end{definition}
\end{definitionbox}

\begin{remark*}[Standard quantified statement]
\[
\begin{aligned}
&\text{Let } E \subseteq \mathbb{R},\ a \in E,\ \text{and } U \subseteq \mathbb{R}.\\
&U \text{ is a neighborhood of } a \text{ in } E
\Longleftrightarrow
\exists \delta > 0\; \bigl(U = E \cap (a-\delta, a+\delta)\bigr).
\end{aligned}
\]
\end{remark*}

\begin{remark*}[Predicate reading]
\[
\operatorname{IsNeighborhood}(U,a,E)
\Longleftrightarrow
\exists \delta > 0\;
\bigl(U = E \cap (a-\delta, a+\delta)\bigr).
\]
\end{remark*}

\begin{remark*}[Negated quantified statement]
\[
\begin{aligned}
&\text{Let } E \subseteq \mathbb{R},\ a \in E,\ \text{and } U \subseteq \mathbb{R}.\\
&U \text{ fails to be a neighborhood of } a \text{ in } E
\Longleftrightarrow
\forall \delta > 0\; \bigl(U \ne E \cap (a-\delta, a+\delta)\bigr).
\end{aligned}
\]
\end{remark*}

\begin{remark*}[Negation predicate reading]
\[
\neg \operatorname{IsNeighborhood}(U,a,E)
\Longleftrightarrow
\forall \delta > 0\;
\bigl(U \ne E \cap (a-\delta, a+\delta)\bigr).
\]
\end{remark*}

\begin{remark*}[Failure modes]
\begin{description}
  \item[Exposition.]
Failure occurs precisely when no radius $\delta$ makes $U$ coincide with the relative interval slice: for every $\delta > 0$ either $U$ contains a point outside $E \cap (a-\delta,a+\delta)$ or the slice contains a point missing from $U$.


  \item[\operatorname{IsNeighborhood}.]
  The displayed obstruction is the corresponding failure mode.
  \[
\begin{aligned}
&\text{Let } E \subseteq \mathbb{R},\ a \in E,\ \text{and } U \subseteq \mathbb{R}.\\
&U \text{ fails to be a neighborhood of } a \text{ in } E
\Longleftrightarrow
\forall \delta > 0\; \bigl(U \ne E \cap (a-\delta, a+\delta)\bigr).
\end{aligned}
  \]
  \[
\neg \operatorname{IsNeighborhood}(U,a,E)
\Longleftrightarrow
\forall \delta > 0\;
\bigl(U \ne E \cap (a-\delta, a+\delta)\bigr).
  \]
\end{description}
\end{remark*}

\begin{remark*}[Failure modes]
\begin{description}
  \item[Exposition.]
\[
\forall \delta > 0\;
\underbrace{\bigl(\exists x \in U:\ x \notin E \cap (a-\delta,a+\delta)\bigr)}_{\text{extraneous point}}
\;\lor\;
\underbrace{\bigl(\exists y \in E \cap (a-\delta,a+\delta):\ y \notin U\bigr)}_{\text{missing point}}.
\]


  \item[\operatorname{IsNeighborhood}.]
  The displayed obstruction is the corresponding failure mode.
  \[
\begin{aligned}
&\text{Let } E \subseteq \mathbb{R},\ a \in E,\ \text{and } U \subseteq \mathbb{R}.\\
&U \text{ fails to be a neighborhood of } a \text{ in } E
\Longleftrightarrow
\forall \delta > 0\; \bigl(U \ne E \cap (a-\delta, a+\delta)\bigr).
\end{aligned}
  \]
  \[
\neg \operatorname{IsNeighborhood}(U,a,E)
\Longleftrightarrow
\forall \delta > 0\;
\bigl(U \ne E \cap (a-\delta, a+\delta)\bigr).
  \]
\end{description}
\end{remark*}

\begin{remark*}[Interpretation]
Relative neighborhoods capture the local behavior of $E$ around $a$ by intersecting $E$ with an ordinary open interval centered at $a$. The construction restricts the ambient notion of openness to the geometry of $E$, so it is suited to relative topologies. Failure simply means that $U$ never matches the precise slice of $E$ at any scale, typically because $U$ omits some nearby points of $E$ or includes points that fall outside $E$ near $a$.
\end{remark*}

\begin{dependencies}
\begin{itemize}
  \item \hyperref[def:subset]{Subset}
  \item \hyperref[def:epsilon-neighbourhood]{Centered Open Interval}
  \item \hyperref[def:order-on-r]{Order on $\mathbb{R}$}
\end{itemize}
\end{dependencies}
% ---------------------------------------------------------
% Definition --- Continuity at a Point
% ---------------------------------------------------------

\begin{definitionbox}{Definition (Continuity at a Point)}
\begin{definition}[Continuity at a Point]
\label{def:continuous-at-point}
Let $A \subseteq \mathbb{R}$, let $f : A \to \mathbb{R}$, and let $c \in A$. The function $f$ is \textbf{continuous at $c$} if and only if
\[
\forall \varepsilon > 0 \;\exists \delta > 0 \;\forall x \in A :
\bigl(|x-c| < \delta \Rightarrow |f(x)-f(c)| < \varepsilon\bigr).
\]
\end{definition}
\end{definitionbox}

\begin{remark*}[Standard quantified statement]
\[
\begin{aligned}
&\text{Let } A \subseteq \mathbb{R},\ f:A\to\mathbb{R},\ c\in A.\\
&f \text{ is continuous at } c
\Longleftrightarrow
\forall \varepsilon>0\;\exists \delta>0\;\forall x\in A\;
\bigl(|x-c|<\delta \Rightarrow |f(x)-f(c)|<\varepsilon\bigr).
\end{aligned}
\]
\end{remark*}

\begin{remark*}[Predicate reading]
\[
  f\colon A\to\mathbb{R},\\

\begin{aligned}
\operatorname{IsContinuous}(f,\mathsf{Point}(c,A),\mathbb R,\mathbb R)
&\Longleftrightarrow
\forall \varepsilon>0\;\exists \delta>0\;\forall x\in A\;
\Bigl[
  \operatorname{Within}(x,I_\delta(c),\mathbb R)
  \Rightarrow
  \operatorname{Within}(f(x),I_\varepsilon(f(c)),\mathbb R)
\Bigr].
\end{aligned}
\]
\end{remark*}

\begin{remark*}[Negated quantified statement]
\[
\begin{aligned}
&\text{Let } A \subseteq \mathbb{R},\ f:A\to\mathbb{R},\ c\in A.\\
&f \text{ is not continuous at } c
\Longleftrightarrow
\exists \varepsilon_0>0\;\forall \delta>0\;\exists x\in A\;
\bigl(|x-c|<\delta \land |f(x)-f(c)| \ge \varepsilon_0\bigr).
\end{aligned}
\]
\end{remark*}

\begin{remark*}[Negation predicate reading]
\[
  f\colon A\to\mathbb{R},\\

\neg\operatorname{IsContinuous}(f,\mathsf{Point}(c,A),\mathbb R,\mathbb R)
\Longleftrightarrow
\exists \varepsilon_0>0\;\forall \delta>0\;\exists x\in A\;
\bigl(|x-c|<\delta \land |f(x)-f(c)| \ge \varepsilon_0\bigr).
\]
\end{remark*}

\begin{remark*}[Failure modes]
\begin{description}
  \item[Exposition.]
Continuity at $c$ fails precisely when there exists a persistent output gap $\varepsilon_0>0$ such that no matter how small the chosen centered interval about $c$ is, a point of $A$ within that interval keeps the function values at least $\varepsilon_0$ away from $f(c)$. The single failure branch records this inescapable oscillation near $c$.


  \item[\operatorname{IsContinuous}.]
  The displayed obstruction is the corresponding failure mode.
  \[
\begin{aligned}
&\text{Let } A \subseteq \mathbb{R},\ f:A\to\mathbb{R},\ c\in A.\\
&f \text{ is not continuous at } c
\Longleftrightarrow
\exists \varepsilon_0>0\;\forall \delta>0\;\exists x\in A\;
\bigl(|x-c|<\delta \land |f(x)-f(c)| \ge \varepsilon_0\bigr).
\end{aligned}
  \]
  \[
\neg\operatorname{IsContinuous}(f,\mathsf{Point}(c,A),\mathbb R,\mathbb R)
\Longleftrightarrow
\exists \varepsilon_0>0\;\forall \delta>0\;\exists x\in A\;
\bigl(|x-c|<\delta \land |f(x)-f(c)| \ge \varepsilon_0\bigr).
  \]
\end{description}
\end{remark*}

\begin{remark*}[Failure modes]
\begin{description}
  \item[Exposition.]
\[
\underbrace{\exists \varepsilon_0>0}_{\text{persistent gap}}
\ \underbrace{\forall \delta>0}_{\text{every centered interval}}
\ \underbrace{\exists x\in A}_{\text{nearby witness}}
:\
|x-c|<\delta \land |f(x)-f(c)| \ge \varepsilon_0.
\]


  \item[\operatorname{IsContinuous}.]
  The displayed obstruction is the corresponding failure mode.
  \[
\begin{aligned}
&\text{Let } A \subseteq \mathbb{R},\ f:A\to\mathbb{R},\ c\in A.\\
&f \text{ is not continuous at } c
\Longleftrightarrow
\exists \varepsilon_0>0\;\forall \delta>0\;\exists x\in A\;
\bigl(|x-c|<\delta \land |f(x)-f(c)| \ge \varepsilon_0\bigr).
\end{aligned}
  \]
  \[
\neg\operatorname{IsContinuous}(f,\mathsf{Point}(c,A),\mathbb R,\mathbb R)
\Longleftrightarrow
\exists \varepsilon_0>0\;\forall \delta>0\;\exists x\in A\;
\bigl(|x-c|<\delta \land |f(x)-f(c)| \ge \varepsilon_0\bigr).
  \]
\end{description}
\end{remark*}

\begin{remark*}[Interpretation]
Continuity at $c$ means that the output $f(x)$ can be kept arbitrarily close to $f(c)$ by taking $x$ sufficiently close to $c$ while staying in the domain $A$. The negated form highlights that a discontinuity arises only when there is a fixed positive separation in output that persists no matter how tightly one zooms in on $c$, so nearby inputs cannot force the outputs to settle near $f(c)$. This captures the geometric intuition that the graph of a continuous function has no abrupt jumps or breaks at $c$.
\end{remark*}

\begin{dependencies}
\begin{itemize}
  \item \hyperref[def:limit-function]{Definition: Limit of a Function}
\end{itemize}
\end{dependencies}

% ---------------------------------------------------------
% Definition --- Continuity (Centered Interval Form)
% ---------------------------------------------------------

\begin{definitionbox}{Definition (Continuity at a Point --- Centered Interval Form)}
\begin{definition}[Continuity at a Point --- Centered Interval Form]
\label{def:continuous-at-point-nbhd}
Let $A \subseteq \mathbb{R}$, let $f : A \to \mathbb{R}$, and let $c \in A$. The function $f$ is continuous at $c$ if and only if for every $\varepsilon > 0$ there exists $\delta > 0$ such that
\[
f\bigl(A \cap I_\delta(c)\bigr) \subseteq I_\varepsilon(f(c)).
\]
\end{definition}
\end{definitionbox}

\begin{remark*}[Standard quantified statement]
\[
\begin{aligned}
&\text{Let } A \subseteq \mathbb{R},\; f:A \to \mathbb{R},\; c \in A.\\
&f \text{ is continuous at } c \Longleftrightarrow \forall \varepsilon > 0\; \exists \delta > 0:\\
&\qquad f\bigl(A \cap I_\delta(c)\bigr) \subseteq I_\varepsilon(f(c)).
\end{aligned}
\]
\end{remark*}

\begin{remark*}[Predicate reading]
\[
  f\colon A\to\mathbb{R},\\

\operatorname{IsContinuous}(f,\mathsf{Point}(c,A),\mathbb R,\mathbb R)
\Longleftrightarrow
\forall \varepsilon > 0\; \exists \delta > 0:\;
f\bigl(A \cap I_\delta(c)\bigr) \subseteq I_\varepsilon(f(c)).
\]
\end{remark*}

\begin{remark*}[Negated quantified statement]
\[
\begin{aligned}
&\text{Let } A \subseteq \mathbb{R},\; f:A \to \mathbb{R},\; c \in A.\\
&\neg\bigl(f \text{ is continuous at } c\bigr) \Longleftrightarrow \exists \varepsilon_{0} > 0\; \forall \delta > 0:\\
&\qquad f\bigl(A \cap I_\delta(c)\bigr) \nsubseteq I_{\varepsilon_0}(f(c)).
\end{aligned}
\]
\end{remark*}

\begin{remark*}[Negation predicate reading]
\[
  f\colon A\to\mathbb{R},\\

\operatorname{IsDiscontinuous}(f,\mathsf{Point}(c,A),\mathbb R,\mathbb R)
\Longleftrightarrow
\exists \varepsilon_{0} > 0\; \forall \delta > 0:\;
f\bigl(A \cap I_\delta(c)\bigr) \nsubseteq I_{\varepsilon_0}(f(c)).
\]
\end{remark*}

\begin{remark*}[Failure modes]
\begin{description}
  \item[Exposition.]
Failure occurs precisely when some fixed output tolerance cannot be matched by any centered interval about $c$. This manifests by finding points of $A$ arbitrarily close to $c$ whose images either exceed the upper endpoint $f(c)+\varepsilon_{0}$ or fall below the lower endpoint $f(c)-\varepsilon_{0}$.


  \item[\operatorname{IsDiscontinuous}.]
  The displayed obstruction is the corresponding failure mode.
  \[
\begin{aligned}
&\text{Let } A \subseteq \mathbb{R},\; f:A \to \mathbb{R},\; c \in A.\\
&\neg\bigl(f \text{ is continuous at } c\bigr) \Longleftrightarrow \exists \varepsilon_{0} > 0\; \forall \delta > 0:\\
&\qquad f\bigl(A \cap I_\delta(c)\bigr) \nsubseteq I_{\varepsilon_0}(f(c)).
\end{aligned}
  \]
  \[
\operatorname{IsDiscontinuous}(f,\mathsf{Point}(c,A),\mathbb R,\mathbb R)
\Longleftrightarrow
\exists \varepsilon_{0} > 0\; \forall \delta > 0:\;
f\bigl(A \cap I_\delta(c)\bigr) \nsubseteq I_{\varepsilon_0}(f(c)).
  \]
\end{description}
\end{remark*}

\begin{remark*}[Failure modes]
\begin{description}
  \item[Exposition.]
\[
\exists \varepsilon_{0} > 0\; \forall \delta > 0\;
\underbrace{\Bigl(
\exists x \in A \cap (c-\delta,c+\delta): f(x) \ge f(c)+\varepsilon_{0}
\Bigr)}_{\text{overshoot}}
\;\lor\;
\underbrace{\Bigl(
\exists x \in A \cap (c-\delta,c+\delta): f(x) \le f(c)-\varepsilon_{0}
\Bigr)}_{\text{undershoot}}.
\]


  \item[\operatorname{IsDiscontinuous}.]
  The displayed obstruction is the corresponding failure mode.
  \[
\begin{aligned}
&\text{Let } A \subseteq \mathbb{R},\; f:A \to \mathbb{R},\; c \in A.\\
&\neg\bigl(f \text{ is continuous at } c\bigr) \Longleftrightarrow \exists \varepsilon_{0} > 0\; \forall \delta > 0:\\
&\qquad f\bigl(A \cap I_\delta(c)\bigr) \nsubseteq I_{\varepsilon_0}(f(c)).
\end{aligned}
  \]
  \[
\operatorname{IsDiscontinuous}(f,\mathsf{Point}(c,A),\mathbb R,\mathbb R)
\Longleftrightarrow
\exists \varepsilon_{0} > 0\; \forall \delta > 0:\;
f\bigl(A \cap I_\delta(c)\bigr) \nsubseteq I_{\varepsilon_0}(f(c)).
  \]
\end{description}
\end{remark*}

\begin{remark*}[Interpretation]
Continuity at $c$ demands that every output tolerance around $f(c)$ has a matching centered interval around $c$ whose image remains confined to the tolerance band. The negation says that one tolerance band resists every centered-interval choice, so some sequence of points from $A$ approaches $c$ while their images either spike above or dip below $f(c)$ by at least a fixed amount. Geometrically, arbitrarily tight vertical strips around the graph above $c$ must be achievable by shrinking the horizontal window; failure means the graph keeps escaping that strip no matter how closely the input variable hugs $c$.
\end{remark*}

\begin{dependencies}
\begin{itemize}
  \item \hyperref[def:continuous-at-point]{Continuity at a Point}
  \item \hyperref[def:relative-neighborhood]{Relative Neighborhood}
  \item \hyperref[def:image-set]{Image of a Set}
  \item \hyperref[def:open-interval]{Open Interval}
\end{itemize}
\end{dependencies}
% ---------------------------------------------------------
% Definition --- Continuity (Sequential)
% ---------------------------------------------------------

\begin{definitionbox}{Definition (Continuity at a Point --- Sequential Form)}
\begin{definition}[Continuity at a Point --- Sequential Form]
\label{def:continuous-at-point-seq}
Let $A \subseteq \mathbb{R}$, let $c \in A$, and let $f \colon A \to \mathbb{R}$. The function $f$ is continuous at $c$ if and only if for every sequence $(x_n)_{n\in\mathbb{N}}$ in $A$, the convergence $\lim_{n\to\infty} x_n = c$ implies $\lim_{n\to\infty} f(x_n) = f(c)$.
\end{definition}
\end{definitionbox}

\begin{remark*}[Standard quantified statement]
\[
\begin{aligned}
&\text{Let } A \subseteq \mathbb{R},\; c \in A,\; f \colon A \to \mathbb{R}.\\
&f \text{ is continuous at } c \Longleftrightarrow
\forall (x_n)_{n\in\mathbb{N}} \subseteq A\;
\Bigl(\lim_{n\to\infty} x_n = c \Rightarrow \lim_{n\to\infty} f(x_n) = f(c)\Bigr).
\end{aligned}
\]
\end{remark*}

\begin{remark*}[Predicate reading]
\[
  f\colon A\to\mathbb{R},\\

\operatorname{IsContinuous}(f,\mathsf{Point}(c,A),\mathbb R,\mathbb R)
\Longleftrightarrow
\forall x_n\;
\Bigl(\mathsf{Sequence}(x_n, A)
\land \operatorname{ConvergesTo}(x_n,c,\mathbb R)
\Rightarrow \operatorname{ConvergesTo}(f(x_n),f(c),\mathbb R)\Bigr).
\]
\end{remark*}

\begin{remark*}[Negated quantified statement]
\[
\begin{aligned}
&\text{Let } A \subseteq \mathbb{R},\; c \in A,\; f \colon A \to \mathbb{R}.\\
&f \text{ is not continuous at } c \Longleftrightarrow
\exists (x_n)_{n\in\mathbb{N}} \subseteq A\;
\Bigl(\lim_{n\to\infty} x_n = c \land \lim_{n\to\infty} f(x_n) \ne f(c)\Bigr).
\end{aligned}
\]
\end{remark*}

\begin{remark*}[Negation predicate reading]
\[
  f\colon A\to\mathbb{R},\\

\neg\operatorname{IsContinuous}(f,\mathsf{Point}(c,A),\mathbb R,\mathbb R)
\Longleftrightarrow
\exists x_n\;
\Bigl(\mathsf{Sequence}(x_n, A)
\land \operatorname{ConvergesTo}(x_n,c,\mathbb R)
\land \neg\operatorname{ConvergesTo}(f(x_n),f(c),\mathbb R)\Bigr).
\]
\end{remark*}

\begin{remark*}[Failure modes]
\begin{description}
  \item[Exposition.]
Continuity at $c$ fails sequentially when there exists a sequence $(x_n)_{n\in\mathbb{N}}$ in $A$ with $x_n \to c$ but $f(x_n)$ does not converge to $f(c)$. This manifests in two distinct ways: either $f(x_n)$ converges to some limit different from $f(c)$, or $f(x_n)$ fails to converge entirely (oscillating without settling or diverging to infinity).


  \item[\operatorname{IsContinuous}.]
  The displayed obstruction is the corresponding failure mode.
  \[
\begin{aligned}
&\text{Let } A \subseteq \mathbb{R},\; c \in A,\; f \colon A \to \mathbb{R}.\\
&f \text{ is not continuous at } c \Longleftrightarrow
\exists (x_n)_{n\in\mathbb{N}} \subseteq A\;
\Bigl(\lim_{n\to\infty} x_n = c \land \lim_{n\to\infty} f(x_n) \ne f(c)\Bigr).
\end{aligned}
  \]
  \[
\neg\operatorname{IsContinuous}(f,\mathsf{Point}(c,A),\mathbb R,\mathbb R)
\Longleftrightarrow
\exists x_n\;
\Bigl(\mathsf{Sequence}(x_n, A)
\land \operatorname{ConvergesTo}(x_n,c,\mathbb R)
\land \neg\operatorname{ConvergesTo}(f(x_n),f(c),\mathbb R)\Bigr).
  \]
\end{description}
\end{remark*}

\begin{remark*}[Failure modes]
\begin{description}
  \item[Exposition.]
\[
\neg(\text{$f$ continuous at $c$}) \Longleftrightarrow
\exists (x_n)_{n\in\mathbb{N}} \subseteq A :\;
\underbrace{\lim_{n\to\infty} x_n = c}_{\text{valid approach}}
\;\land\;
\underbrace{\lim_{n\to\infty} f(x_n) \ne f(c)}_{\text{image sequence fails}}.
\]


  \item[\operatorname{IsContinuous}.]
  The displayed obstruction is the corresponding failure mode.
  \[
\begin{aligned}
&\text{Let } A \subseteq \mathbb{R},\; c \in A,\; f \colon A \to \mathbb{R}.\\
&f \text{ is not continuous at } c \Longleftrightarrow
\exists (x_n)_{n\in\mathbb{N}} \subseteq A\;
\Bigl(\lim_{n\to\infty} x_n = c \land \lim_{n\to\infty} f(x_n) \ne f(c)\Bigr).
\end{aligned}
  \]
  \[
\neg\operatorname{IsContinuous}(f,\mathsf{Point}(c,A),\mathbb R,\mathbb R)
\Longleftrightarrow
\exists x_n\;
\Bigl(\mathsf{Sequence}(x_n, A)
\land \operatorname{ConvergesTo}(x_n,c,\mathbb R)
\land \neg\operatorname{ConvergesTo}(f(x_n),f(c),\mathbb R)\Bigr).
  \]
\end{description}
\end{remark*}

\begin{remark*}[Interpretation]
The sequential characterization asserts that information about $f$ near $c$ is completely encoded by its action on limits of sequences taken within the domain. Any approach to $c$ through admissible points must transport convergent input sequences to convergent output sequences with the expected limit; conversely, a single pathological sequence suffices to expose a discontinuity. Geometrically, this captures the idea that $f$ cannot introduce jumps or oscillations that persist along every infinitesimal path toward $c$.
\end{remark*}

\begin{dependencies}
\hyperref[def:continuous-at-point]{Definition~\ref*{def:continuous-at-point}}.
\end{dependencies}
